\subsection{基环扩张}

设 \(\mathbf{K}_1\) 是带单位元的交换环，\(\phi\) 是从 K 到 \(\mathbf{K}_1\) 的把 1 映为 1 的同态。设 \(\mathfrak{g}\) 是 李 K-代数，U 是其包络代数，M 是左 g-模，即左 U-模。于是 \(\mathrm{M}_{(\mathbf{K}_1)}\) 具有典范的左 \(\mathrm{U}_{(\mathbf{K}_1)}\)-模结构，因而具有左 \(\mathfrak{g}_{(\mathbf{K}_1)}\)-模结构。设 \(\rho\)、\(\rho_{(\mathbf{K}_1)}\) 分别是对应于 M 和 \(\mathfrak{g}\) 的 \(\mathfrak{g}_{(\mathbf{K}_1)}\) 和 \(\mathrm{M}_{(\mathbf{K}_1)}\) 的表示；称 \(\rho_{(\mathbf{K}_1)}\) 是由 \(\rho\) 通过扩张基环导出的表示，并可应用《代数》第 VIII 章 § 13 第 4 号的结果。若 \(x \in \mathfrak{g}\)，则 \(\rho_{(\mathbf{K}_1)}(x)\) 正是 \(\rho(x) \otimes 1\) 上的自同态 \(\mathrm{M}_{(\mathbf{K}_1)} = \mathrm{M} \otimes_{\mathbf{K}} \mathbf{K}_1\)。

假设 K 是域，\(K_{1}\) 是 K 的扩域，\(\phi\) 是 K 到 \(K_{1}\) 的典范嵌入。设 V 和 \(V'\) 是 M 的向量子空间。令 a 为由满足 \(x \in g\) 的 \(\rho(x)(V) \subset V'\) 组成的 g 的向量子空间。令 \(a'\) 为由满足 \(\mathfrak{g}_{(\mathbf{K}_{1})}\) 的 \(x' \in \mathfrak{g}_{(\mathbf{K}_{1})}\) 组成的 \(\rho_{(\mathbf{K}_{1})}(x')(V_{(\mathbf{K}_{1})}) \subset V'_{(\mathbf{K}_{1})}\) 的向量子空间。则 \(a' = a_{(\mathbf{K}_{1})}\)。显然 \(a_{(\mathbf{K}_{1})} \subset a'\)。现在令 \(x' \in a'\)。可写成 \(x' = \sum_{i=1}^{n} \lambda_{i} x_{i}\)，其中 \(x_{i}\) 属于 g，\(\lambda_{i}\) 是 \(K_{1}\) 中在 K 上线性无关的元素。对所有 \(u \in V\)，\(\rho(x') .u \in V'_{(\mathbf{K}_{1})}\)，即 \(\sum_{i=1}^{n} \lambda_{i} \rho(x_{i}) .u \in V'_{(\mathbf{K}_{1})}\)，从而 \(\rho(x_{i}) .u \in V'\)，于是 \(x_{i} \in a\) 且 \(x' \in a_{(\mathbf{K}_{1})}\)。这说明 \(a' = a_{(\mathbf{K}_{1})}\)。特别地，\(\mathfrak{g}_{(\mathbf{K}_{1})}\) 的中心由扩张 K 到 \(K_{1}\) 时 g 的中心导出：只需将上面的结论应用于 g 的伴随表示即可。因此，对所有 p.~都有 \(\mathcal{C}_{\rho'}(\mathfrak{g}_{(\mathbf{K}_{1})}) = (\mathcal{C}_{\rho}\mathfrak{g})_{(\mathbf{K}_{1})}\)。类似地，设 h 是 g 的子代数，n 是 h 在 g 中的正规化子。则 \(\mathfrak{h}_{(\mathbf{K}_{1})}\) 在 \(\mathfrak{g}_{(\mathbf{K}_{1})}\) 中的正规化子是 \(n_{(\mathbf{K}_{1})}\)。

设 \(\mathbf{K}, \mathbf{K}_{1}, \mathfrak{g}, \rho, \mathbf{M}\) 如上一段所述。设 \(\mathfrak{b}\) 是 \(\mathfrak{g}\) 的向量子空间，\(\mathbf{W}\) 是 \(\mathbf{M}\) 的向量子空间。令 \(\mathbf{V}\) 为由满足 \(\mathbf{M}\) 的 \(m \in \mathbf{M}\) 组成的 \(\rho(\mathfrak{b}). m \subset W\) 的向量子空间。令 \(V'\) 为由满足 \(\mathrm{M}_{(\mathbf{K}_{1})}\) 的 \(m' \in \mathrm{M}_{(\mathbf{K}_{1})}\) 组成的 \(\rho_{(\mathbf{K}_{1})}(\mathfrak{b}_{(\mathbf{K}_{1})}). m' \subset W_{(\mathbf{K}_{1})}\) 的向量子空间。如上可见，\(V' = V_{(K_1)}\)。特别地，\(M_{(K_1)}\) 的不变元向量子空间由 M 的不变元向量子空间通过将基域从 K 扩张到 \(K_1\) 而得到。

令 \(K, K_{1}\) 和 \(\phi\) 如本号开头所述。设 g 是 李 K-代数，M、N 是 g-模。若 M 和 N 是同构的 g-模，则 \(M_{(K_{1})}\) 和 \(N_{(K_{1})}\) 是同构的 \(g_{(K_{1})}\)-模。反之：

命题 13. 设 K 是域，\(K_{1}\) 是 K 的扩域，g 是 李 K-代数，M、N 是在 K 上有限维的两个 g-模。若 \(\mathrm{M}_{(\mathbf{K}_{1})}\) 和 \(\mathrm{N}_{(\mathbf{K}_{1})}\) 是同构的 \(\mathfrak{g}_{(\mathbf{K}_{1})}\)-模，则 M、N 是同构的 g-模。

证明分两步。

\begin{enumerate}
\def\labelenumi{(\arabic{enumi})}
\item
  首先假设 \(\mathbf{K}_1\) 是 K 的有限扩域，次数为 \(n\)。设 U 是 \(\mathfrak{g}\) 的包络代数，于是 \(\mathfrak{g}_{(\mathbf{K}_1)}\) 的包络代数是 \(\mathrm{U}_{(\mathbf{K}_1)} = \mathrm{U} \otimes_{\mathbb{K}} \mathrm{K}_1\)（§ 2，第 9 号）。由于 \(\mathrm{M}_{(\mathbf{K}_1)}\) 和 \(\mathrm{N}_{(\mathbf{K}_1)}\) 作为 \(\mathrm{U}_{(\mathbf{K}_1)}\)-模同构，它们当然作为 U-模也同构；但作为 U-模，它们分别同构于 \(\mathbf{M}^n\) 和 \(\mathbf{N}^n\)。现在 M、N 是有限长度的 U-模；因此 M（分别 N）是子模族 \((\mathrm{P}_i^{s_i})_{1 \leqslant i \leqslant p}\)（分别 \((\mathrm{Q}_j^{s_j})_{1 \leqslant j \leqslant q}\)）的直和，其中 \(\mathrm{P}_i\)（分别 \(\mathrm{Q}_j\)）不可分解，并且不同指标的两个 \(\mathrm{P}_i\)（分别 \(\mathrm{Q}_j\)）不同构（《代数》，第 VIII 章，§ 2，第 2 号定理 1）。于是 \(\mathrm{M}^n\)（分别 \(\mathrm{N}^n\)）同构于 \(\mathrm{P}_i^{nr_i}\)（分别 \(\mathrm{Q}_j^{ns_i}\)）的直和；由此（同上）得 \(p = q\)，并且必要时置换 \(\mathrm{Q}_j\) 后，\(nr_i = ns_i\)，且 \(\mathrm{P}_i\) 时 \(\mathrm{Q}_i\) 同构于 \(1 \leqslant i \leqslant p\)，所以 M 同构于 N。
\item
  一般情形。令 P 为 g-模 \(\mathcal{L}_{\mathrm{K}}(\mathrm{M},\mathrm{N})\)，Q 为 P 的不变元子空间，即从 g-模 M 到 g-模 N 的同态组成的集合。在 \(\mathfrak{g}_{(\mathrm{K}_{1})}\)-模 \(\mathcal{L}_{\mathrm{K}_{1}}(\mathrm{M}_{(\mathrm{K}_{1})},\mathrm{N}_{(\mathrm{K}_{1})})=(\mathcal{L}_{\mathrm{K}}(\mathrm{M},\mathrm{N})_{\mathrm{K}_{1}})\) 中，不变元子空间是 \(\mathrm{Q}_{(\mathrm{K}_{1})}\)。假设 \(\mathrm{M}_{(\mathrm{K}_{1})}\) 和 \(\mathrm{N}_{(\mathrm{K}_{1})}\) 同构，则 M 和 N 在 K 上有相同的维数，并且 \(\mathrm{Q}_{(\mathrm{K}_{1})}\) 中存在元素 g，它是从 \(\mathrm{M}_{(\mathrm{K}_{1})}\) 到 \(\mathrm{N}_{(\mathrm{K}_{1})}\) 的同构。令 \((f_{1},\ldots,f_{d})\) 是 Q 在 K 上的基，并选取 M、N 在 K 上的基。若
\end{enumerate}

\[
f = \sum_ {k = 1} ^ {d} \lambda_ {k} f _ {k}
\]

若 \(\lambda_{k}\in K_{1}\) 且 \(1\leqslant k\leqslant d\)，则相对于这些基，\(f=\sum_{k=1}\lambda_{k}f_{k}\) 的矩阵的行列式是一个系数在 K 中的多项式 \(D(\lambda_{1},\ldots,\lambda_{d})\)。当 f=g 时，该行列式非零，因而 D 的系数不全为非零。因此，若 \(\Omega\) 是 K 的代数闭包，则存在（因为 \(\Omega\) 是无限的）元素 \(\mu_{k}\in\Omega\)（\(1\leqslant k\leqslant d\)），使得 \(D(\mu_{1},\ldots,\mu_{d})\neq0\)（《代数》，第 IV 章，§ 2，第 5 号命题 8）。若 \(K_{2}\) 是由

\[
\sum_ {k = 1} ^ {d} \mu_ {k} f _ {k}
\]

是由 \(\mu_{k}\)（\(1 \leqslant k \leqslant d\)）生成的 K 的代数扩域，则 \(\sum_{k=1}^{d} \mu_{k} f_{k}\) 是从 \(\mathrm{M}_{(\mathbf{K}_{2})}\) 到 \(\mathrm{N}_{(\mathbf{K}_{2})}\) 的同构；但 \(\mathbf{K}_{2}\) 在 K 上是有限次的（《代数》，第 V 章，§ 3，第 2 号命题 5），所以由论证的第一部分，M 和 N 同构。

再次令 \(K, K_{1}\) 和 \(\phi\) 如本号开头所述。设 \(\rho\) 是 \(\mathfrak{g}\) 在有有限基 \((x_{1},\ldots ,x_{n})\) 的 K-模 M 上的表示。则 \(\mathfrak{g}_{(\mathrm{K}_1)}\) 上与 \(\rho_{(\mathrm{K}_1)}\) 关联的双线性形式，是将与 \(\rho\) 关联的双线性形式的基环扩张到 \(\mathrm{K}_1\) 后得到的（因为若 \(u\in \mathcal{L}_{\mathrm{K}}(\mathbf{M})\)，\(u\) 相对于 \((x_{1},\dots,x_{n})\) 的矩阵，与 \(u\otimes 1\) 相对于 \((x_{1}\otimes 1,\dots,x_{n}\otimes 1)\) 的矩阵相同，从而 \(u\) 与 \(u\otimes 1\) 有相同的迹）。特别地，若 K-模 \(\mathfrak{g}\) 有有限基，则 \(\mathfrak{g}_{(\mathrm{K}_1)}\) 的 Killing 形式由 \(\mathfrak{g}\) 的 Killing 形式通过将基环扩张到 \(\mathrm{K}_1\) 而得到。

